\section{Conclusion}
In this work, we present \our{}, an orchestration-centric agentic system that automates sub-agent creation through a unified four-tuple interface (Instruction, Context, Tools, Model) to solve complex, long-horizon agentic tasks. By treating sub-agents as dynamically creatable units, the orchestrator can spawn task-tailored executors on demand with specialized working memory (instruction, context) and capabilities (model, tools). 

This abstraction brings practical benefits: it enables on-demand specialization with task-sufficient context, keeps sub-agents plug-and-play across implementations.
The decoupling of abstraction from execution makes it learnable, and we present two ways to optimize the orchestrator through supervised fine-tuning and context learning. 
Empirically, \our{} demonstrates strong and consistent improvements across three challenging benchmarks (GAIA, Terminal-Bench, and SWE-Bench-Verified) when paired with frontier models like Gemini-3-Flash. It significantly outperforms established baseline frameworks, achieving, for instance, an average gain of 16.28\% points in pass@1 across all benchmarks. These results validate the effectiveness of our orchestration-centric approach to automating complex, long-horizon tasks.